\chapter{The Full Battery on Scholar's A40, and Stage 6 for the First Time}
\label{chap:s2-1}

\section{Goal}

Every Study~2 number so far came from one RTX~3070 with 8~GB of memory, where a single Stage~2 pass takes
about ten seconds a question. Stage~6 (stability: five sampled reruns per row) was implemented and tested but
never run, because at that speed it would have occupied the card for days. This phase moves the whole
battery to Purdue's Scholar cluster (NVIDIA A40, 46~GB), runs every stage there including Stage~6, and
checks that the headline figures survive the change of machine.

\section{What Was Done}

\textbf{Per-machine result trees.} A parity run had already shown that greedy decoding reproduces exactly on
one GPU but not across GPUs: on identical weights, code, data and pinned libraries the A40 agreed with the
frozen 3070 baseline on 96.9\% of verdicts and 42.7\% of replies byte-for-byte. Every model stage therefore
records the GPU in its fingerprint and refuses to resume a file written on another device, and the Scholar
run writes to its own tree, \texttt{results/scholar-a40/}. Comparisons between stages (Stage~6 against
Stage~2, for instance) stay inside one tree, so machine noise is never measured as instability.

\textbf{Self-continuing job chains.} Scholar caps a job at four hours and lets one user run two at once.
Each link of a chain queues its own successor first (\texttt{afterany}), runs its stages under
\texttt{timeout -{}-signal=INT} about 25 minutes short of the wall so a stage stops between rows rather than
mid-write, and on success marks itself done and cancels the successor. A real failure writes a
\texttt{FAILED} marker and stops every chain for that model --- there is no retry loop. Four chains
(Stages 4, 5, 6 and later 8) run after the base chain (Stages 0--3, 7), and the last to finish submits the
final chain (Stages 9--11).

\section{Results}

\begin{center}
\begin{tabular}{p{6.2cm}rr}
\toprule
\textbf{metric} & \textbf{RTX 3070 (frozen)} & \textbf{A40} \\
\midrule
sensitivity                         & 90.9\% & 90.4\% \\
specificity                         & 40.9\% & 40.5\% \\
false-alarm rate                    & 59.1\% & 59.5\% \\
rule\_4 sensitivity ($n=70$)        & 30.0\% & 25.7\% \\
Stage 4 flip rate (black box)       & 6.4\%  & 6.4\%  \\
Stage 5 attention exceeds area null & 85.7\% & 85.4\% \\
Stage 10 explanation correct        & 2.4\%  & 2.4\%  \\
\bottomrule
\end{tabular}
\captionof{table}{The same battery on two machines. Every headline agrees within about one point and well
inside its 95\% interval; rule\_4's four-point gap is three images out of 70.}
\label{tab:s21-parity}
\end{center}

\textbf{Stage 6, run for the first time} (799 rows: 400 baseline-YES, 399 baseline-NO; five reruns each at
temperature 0.7):

\begin{center}
\begin{tabular}{p{7.4cm}r}
\toprule
\textbf{metric} & \textbf{A40, 4-bit} \\
\midrule
verdict unanimous across all five reruns    & 51.9\% [48.4, 55.4] \\
\quad baseline YES / baseline NO            & 54.2\% / 49.6\% \\
majority of reruns overturns the greedy verdict & 17.9\% [15.3, 20.7] \\
flickering region (present in some reruns, not all) & 38.0\% [34.7, 41.4] \\
region drift $>$ 10\% of the image diagonal ($n=612$) & 36.1\% [32.4, 40.0] \\
\bottomrule
\end{tabular}
\captionof{table}{Stage 6 stability. Stage 6's own console summary printed a median pairwise agreement of
100\%, which reads as stable; unanimity shows half the rows disagree with themselves at least once.}
\label{tab:s21-stability}
\end{center}

On nearly one row in five, the verdict the paper reports (greedy decoding) is the \emph{minority} answer
under the model's own sampling distribution. The explanation is less stable still: where at least two
reruns named a region, a third of the time those regions sit more than a tenth of the image diagonal apart.

\section{Problems and Solutions}

\textbf{Stage 0's small-object score dropped from 89\% to 78\%.} That is 8 of 9 images against 7 of 9 --- one
image. The verdict becomes ``GO with limitation'' on the A40, and the paper should not read anything into a
one-image difference beyond that.

\textbf{The ``3.5$\times$ unexplained latency gap'' was the local machine.} On the 3070 the same images
answered in 3.1~s standalone and 10--11~s inside a run, and the latency caveat in Stage~7 and the runner
blamed the harness. On the A40 the in-run median is 2.3~s. The gap belongs to the 3070 machine; the caveat
now quotes a standalone-to-in-run ratio only when both were measured on the same GPU.

\textbf{Region drift was never computed.} Stage~6 reads its rows from Stage~2, which carries no image size,
so the drift column stayed empty. Stage~3's regions file has the size of every (image, rule), so the report
now joins on it and recomputes drift from the stored regions; future Stage~6 runs take the sizes in-run.
No other figure moved.

\section{Standard vs.\ Non-Standard Choices}

\textbf{Standard.} Bootstrap intervals on every rate; resumable, checkpointed runs.

\textbf{Non-standard.} Treating the GPU as part of the experimental configuration --- fingerprinted, never
mixed within one tree --- is stricter than common practice, and it is what the 96.9\% parity figure demands:
3.1\% of greedy verdicts change with the card alone.

\section{Implications}

The headline results are not an artefact of one machine. Stage~6 adds a new finding for the paper: the
greedy verdict is the minority answer in 17.9\% of rows, and explanation regions are less stable than the
verdicts they accompany.
